\section{Audit Semantics}

The verifier contract is profile-scoped. It can confirm that declared
constraints, lineage, canonical ordering, hashes, and package contents match
the profile. It cannot certify fairness, neutrality, or legal sufficiency
outside the declared audit profile and available context.

Solver reports, heuristic transcripts, and selected-frontier manifests are
upstream evidence. The certificate consumes their final exported artifacts and
checks the package under a declared profile. This keeps exact proofs,
heuristic-improvement summaries, and Pareto-selection records from being
collapsed into one ambiguous claim.

A successful verification should therefore be read as: this plan, with this
context, under this profile, with these hashes and lineage records, passed the
declared checks. It should not be read as: this plan is fair, legally adequate,
or preferred among all available alternatives.

The public verifier workflow is intentionally fixed point shaped: algorithm
output becomes a package, the package certificate is checked, and the checked
package can then be referenced by reports. The neutral surface is
the RPLAN certificate verifier; the BISECT manifest verifier consumes the same
package manifest shape as a bridge.
